\documentclass[14pt,a4paper]{extarticle}
\usepackage[margin=20mm]{geometry}
\usepackage{fontspec}
\setmainfont{Times New Roman}
\setsansfont{Arial}
\setmonofont{Consolas}[Scale=MatchLowercase]
\usepackage{polyglossia}
\setdefaultlanguage{russian}
\setotherlanguage{english}
\usepackage{amsmath,booktabs,array,graphicx,needspace}
\usepackage[hyphens]{url}
\usepackage[hidelinks]{hyperref}
\setlength{\parindent}{1.25cm}
\setlength{\parskip}{0pt}
\setlength{\emergencystretch}{2em}
\clubpenalty=10000
\widowpenalty=10000
\sloppy
\hypersetup{pdftitle={Ресурсная эффективность локальных языковых моделей в аграрных задачах: сравнение конфигураций 9B и 27B и использование расчётных инструментов},pdfauthor={Антоний Сарычев; Денис Шестернин},pdfsubject={Working empirical draft; final author consent and license pending}}
\begin{document}
\begin{center}\bfseries Ресурсная эффективность локальных языковых моделей в аграрных задачах: сравнение конфигураций 9B и 27B и использование расчётных инструментов\end{center}

Антоний Сарычев, Денис Шестернин


РГАУ-МСХА имени К.А. Тимирязева, Москва, Россия


Рабочий черновик от 4 октября 2026 года. RU/EN описывают одно исследование; окончательное согласие авторов и лицензия ожидаются.


УДК 004.8 --- рабочий индекс, требует подтверждения выбранной площадкой.


\section*{Аннотация}

Сопоставлены качество, задержка и память локальных конфигураций Qwen 9B и 27B на аграрном DEV-наборе. Проверена совместимость исторических ответов и телеметрии. Расчётные задачи рассматриваются отдельно от знаний; декларация инструмента отделена от его исполнения.


Ключевые слова: языковые модели; аграрные задачи; локальный вывод; ресурсы; детерминированные расчёты; воспроизводимость.


\section*{1. Проблема и цель исследования}

Локальному аграрному ассистенту могут поступать вопросы на предметные знания и задачи с явно заданной формулой. Связный ответ не гарантирует верного расчёта, а название правильного инструмента ещё не означает его вызов. При ограниченной памяти домашнего компьютера выбор модели должен учитывать качество всей конфигурации, задержку, занятую память и наблюдавшиеся отказы.


Цель работы --- описать практический компромисс между Qwen3.5-9B и Qwen3.8-27B и проверить роль реального исполнения расчётного кода. Основные вопросы: какое качество и ресурсы наблюдаются у этих конфигураций; как отличается самостоятельный расчёт от передачи аргументов инструменту. Облачное делегирование остаётся дополнительным вопросом. Сравнение меняет поколение модели, размер, квантование и размещение в памяти; причинный эффект числа параметров отдельно не выделяется.


Основной материал --- публичный DEV MERA ruTXTAgroBench. Отдельный TELEMETRY40 связывает сохранённые ответы с измерениями ресурсов, а Study C проверяет заданные формулы, единицы и недостающие входы. Исторический HOME150 служит дополнительным контекстом. Точности этих блоков не объединяются. Независимый скрытый TEST в этой сессии не запускался. 9B --- фокусная практическая конфигурация, 27B --- reference с более высоким наблюдаемым качеством.


\section*{2. Связанные исследования}

FrugalGPT, AutoMix, Hybrid LLM, RouteLLM и единая теория маршрутизации рассматривают выбор модели и эскалацию \cite{ref1,ref2,ref3,ref4,ref5}. RouterBench показывает значение различия между анализом сохранённых ответов и работающим маршрутизатором \cite{ref6}. Их опубликованные преимущества не переносятся автоматически на нашу машину и набор. Работы об оценке неопределённости и отказе от ответа требуют отдельного критерия надёжности \cite{ref7,ref8}; валидный JSON или самооценка уверенности таким критерием не являются.


PLLaMa, AgriGPT и AgroLLM исследуют аграрную адаптацию и передачу знаний \cite{ref9,ref10,ref11}. AgriAgent и AgroTools уже рассматривают планирование и инструменты в сельском хозяйстве \cite{ref12,ref13}. Поэтому вклад настоящей работы ограничен практическим сравнением локальных ресурсов и учётом завершённого расчёта. tau-bench оценивает достигнутый результат взаимодействия, ToolACE --- данные и проверку вызовов функций \cite{ref14,ref15}.


BitNet и технический отчёт Gemma показывают разные подходы к эффективности и архитектурам \cite{ref16,ref17}, но не дают результатов для наших GGUF. AgXQA относится к другому виду аграрных вопросов, использованному в историческом HOME \cite{ref18}. Для вычислений отдельно применяются фундаментальные формулы FAO-56 \cite{ref19}. Карты моделей, API и программные репозитории учитываются как технические источники и не увеличивают число научных публикаций.


\section*{3. Данные и методика}

ruTXTAgroBench содержит 2 935 русских вопросов по десяти аграрным дисциплинам, включая 293 публичных DEV. Авторы набора сообщают о составлении агрономами, редактировании и проверке тремя экспертами. Это описание создателей, а не наш независимый аудит каждого эталона или доказательство пригодности совета для хозяйства. Вопросы, ключи и полные журналы хранятся локально с ограничениями распространения.


Начальный сбалансированный DEV60 входит в DEV293 и использован при выборе моделей. TELEMETRY40 берёт по четыре позиции в каждой дисциплине из того же DEV без отбора по правильности; независимой проверкой он не является. С первоначальным DEV60 пересекаются 20 его заданий. Study C содержит 27 сконструированных русских заданий с повторяющимися семействами формул. HOME150 --- отдельный исторический контролируемый опыт с синтетическими и двуязычными случаями.


Строгая EM проверяет равенство множеств латинских обозначений выбранных вариантов. F1 --- средняя по заданиям гармоническая мера точности и полноты множества. Четыре исходных эталона смешивают кириллицу и латиницу и не разбираются по заявленному правилу. Их не транслитерировали и не исправляли: основной знаменатель равен 289, все 293 попытки сохранены. Консервативный расчёт с нулём для этих четырёх строк приведён отдельно.


Опубликованный scorer MERA использует SQuAD-нормализацию, которая, в частности, может удалить латинскую A как английский артикль. Его диагностика и наши строгие EM/F1 различаются. Серверная реализация не подтверждена; официального результата MERA мы не заявляем. Интервалы на повторно использованном DEV описательные; допустимая граница потери качества 9B не задана.


Измерения относятся к Ryzen 7 8700F, RTX 5060 Ti с 16 311 MiB VRAM и 31,620 GiB RAM. Зафиксированы Windows 11 Pro, драйвер NVIDIA 617.14, LM Studio 0.4.25+1 и выбранный llama.cpp CUDA12 runtime 2.51.0. Точная сборка движка для каждой модели не раскрыта. Основные GGUF имеют 8,873 GiB для Qwen9 Q8\_0 и 15,334 GiB для Qwen27 UD-Q4\_K\_M; их идентичность определяется квитанциями SHA-256.


Исторический DEV использовал OpenAI-совместимый транспорт, строгую JSON schema, seed 42, контекст 2 048, temperature 0, top\_p 1, максимум 512 токенов и отключённые рассуждения. Запрошенный offload --- max и 0,75, checkpoints --- 32; фактические GPU layers неизвестны. Native TELEMETRY40 сохраняет тексты и основные ограничения декодирования, но не задаёт раскрытый seed или серверную schema. Checkpoints равны нулю, размещение GPU автоматическое, MTP выключен. Последовательность блоков ABBA, отдельный синтетический прогрев и выгрузки раскрыты. Загрузку и прогрев не включают в задержку запроса.


Память измеряется как sampled whole-GPU, whole-system RAM и RSS группы процессов LM Studio. Фоновые приложения и shared pages входят в эти области; выделенная память модели не установлена. Задержка --- клиентский HTTP round trip; клиентская и серверная TTFT различаются. Скорость короткого ответа буквами не характеризует произвольную длинную генерацию. Энергия не измерялась.


\section*{4. Качество и ресурсы}

\noindent\begin{minipage}{\linewidth}\fontsize{12}{14}\selectfont
Таблица 1. Исторический DEV293 в OpenAI-совместимом профиле. Выполнены все 293 запроса каждой модели, доступны 289 разбираемых эталонов.


\begin{center}\begingroup\fontsize{12}{14}\selectfont\setlength{\tabcolsep}{4pt}
\begin{tabular}{@{}>{\raggedright\arraybackslash}p{0.31\linewidth}>{\raggedright\arraybackslash}p{0.24\linewidth}>{\raggedright\arraybackslash}p{0.12\linewidth}>{\raggedright\arraybackslash}p{0.23\linewidth}@{}}
\toprule
Конфигурация & Строгая EM & Set F1 & p50 / p95, с \\
\midrule
Qwen3.5-9B Q8\_0 & 141/289; 48,79\% & 0,6286 & 0,295 / 0,402 \\
Qwen3.8-27B UD-Q4\_K\_M & 150/289; 51,90\% & 0,6468 & 1,374 / 1,753 \\
\bottomrule\end{tabular}\endgroup\end{center}\end{minipage}

Разность Qwen9-Qwen27 составляет -3,11 п.п. Парные исходы: обе верны 113, только 9B --- 28, только 27B --- 37, обе неверны 111. Сохранённый bootstrap по дисциплинам, 10 000 повторов, даёт описательный 95\%-й интервал разности EM [-8,30; +2,42] п.п. и F1 [-0,06148; +0,02552]. Отношение 141/150 = 94\% --- относительное отношение EM; точность Qwen9 равна 48,79\%, эквивалентность не доказана. Консервативная чувствительность равна 141/293 = 48,12\% и 150/293 = 51,19\%; истинные ответы четырёх спорных строк она не устанавливает.


\noindent\begin{minipage}{\linewidth}\fontsize{12}{14}\selectfont
Таблица 2. Native TELEMETRY40: качество сохранённых ответов оценено ретроспективно на 39 эталонах; задержка и ресурсы относятся к 40 завершённым запросам каждой модели.


\begin{center}\begingroup\fontsize{12}{14}\selectfont\setlength{\tabcolsep}{4pt}
\begin{tabular}{@{}>{\raggedright\arraybackslash}p{0.44\linewidth}>{\raggedright\arraybackslash}p{0.23\linewidth}>{\raggedright\arraybackslash}p{0.23\linewidth}@{}}
\toprule
Измерение & Qwen9 & Qwen27 \\
\midrule
Строгая EM & 22/39; 56,41\% & 24/39; 61,54\% \\
Set F1 & 0,7197 & 0,7801 \\
p50 / p95, с & 0,269 / 0,344 & 1,434 / 1,829 \\
Client TTFT p50, с & 0,122 & 0,691 \\
GPU request peak, MiB & 11 471 & 13 287 \\
RAM request peak, GiB & 20,858 & 28,470 \\
RSS request peak, MiB & 11 198,96 & 18 742,34 \\
Завершено / нарушения JSON & 40/40 / 2 & 40/40 / 0 \\
\bottomrule\end{tabular}\endgroup\end{center}\end{minipage}

В native-профиле медиана задержки Qwen9 меньше в 5,338 раза, а Qwen27 верно отвечает на два задания больше. Ни одна конфигурация не лучше по всем критериям. Это описательный компромисс без выбранного после наблюдения порога допустимой потери.


Сверка 80 пар подтверждает одинаковые байты system/user и записи идентичности весов, но выбранные варианты меняются в 1/40 случае Qwen9 и 2/40 Qwen27. Исторические ответы на эти же задания дали обеим моделям 23/39, F1 0,7282 и 0,7374. Поэтому качество DEV293 и новую скорость нельзя выдавать за одновременные измерения при одинаковых настройках. Для TELEMETRY40 совместное описание допустимо с пометкой о ретроспективном scorer и повторном DEV.


Первый native-профиль с 32 checkpoints привёл к остановкам Qwen27 по RAM guard. При checkpoints 0 TELEMETRY40 завершён; отказы и смена конфигурации сохранены. Операционный отказ не считается арифметической ошибкой. Одна успешная серия не устанавливает долговременную устойчивость.


\section*{5. Расчёты и вызовы инструментов}

Во всех ветвях Study C явно доступны тип задачи, вопрос, входные числа и единицы. Проверяются простая суточная сумма тепла без верхнего ограничения, давление насыщения/VPD, суточная эталонная ET0 FAO-56 и заданная экспериментальная разность max(0, target - sum(rain)). GDD суммирует max(0, (Tmin + Tmax)/2 - base) по заданным дням. VPD по одному наблюдению равен давлению насыщения, умноженному на (1 - RH/100); суточная версия использует экстремумы температуры и влажности. Разность осадков не является готовой нормой полива. ET0 отличается от ETc и дозы орошения \cite{ref19}. Недостающие измерения, ветер, G и пороги культуры не подставляются из скрытого эталона или типовых значений.


Историческая полная правильность Qwen9/Qwen27 равна 7/27 и 15/27; число с единицами --- 1/16 и 6/16; правильное название инструмента --- 17/26 и 25/26. Неподдерживаемый случай исключён только из tool-знаменателя. Исполнитель совпал со спецификацией и независимо построенными формульными эталонами в 27/27. Replay выбранных названий дал 17/26 и 25/26, поскольку код получал исходные входы напрямую; передачи аргументов моделью такой опыт не проверял. Исторические ответы 9B и 27B получены при разных checkpoint-профилях, что ограничивает сравнение задержек.


Новый заранее объявленный опыт завершил все 108 решений: 140 измеряемых модельных вызовов и два отдельных синтетических прогрева. Обе модели используют native-профиль checkpoints 0. Ветка без tools делает один вызов; инструментальная ветка выбирает разрешённую функцию и передаёт полные аргументы. Программа проверяет их равенство публичным входам и единицам, затем исполняет формулу и передаёт результат модели для итогового ответа. Из gold, скрытой категории или ожидаемого tool аргументы не дополняются. Это живое JSON-исполнение на уровне приложения с явно заданными типом и контрактами; native function-calling и обнаружение намерения из свободной речи не проверялись.


\noindent\begin{minipage}{\linewidth}\fontsize{12}{14}\selectfont
Таблица 3. Study C, 27 случаев в каждой ветви. Полная правильность сохраняет scorer значений/action/tool; контракт итогового JSON проверен отдельно. Задержка --- сумма модельных HTTP-интервалов, без dispatch и журналирования.


\begin{center}\begingroup\fontsize{12}{14}\selectfont\setlength{\tabcolsep}{4pt}
\begin{tabular}{@{}>{\raggedright\arraybackslash}p{0.28\linewidth}>{\raggedright\arraybackslash}p{0.16\linewidth}>{\raggedright\arraybackslash}p{0.16\linewidth}>{\raggedright\arraybackslash}p{0.16\linewidth}>{\raggedright\arraybackslash}p{0.14\linewidth}@{}}
\toprule
Конфигурация & Верно & Число + единица & Исполнений & p50 суммы, с \\
\midrule
Qwen9 без tools & 7/27 & 1/16 & 0 & 1,469 \\
Qwen9 live tools & 18/27 & 16/16 & 16 & 2,857 \\
Qwen27 без tools & 15/27 & 6/16 & 0 & 8,989 \\
Qwen27 live tools & 25/27 & 16/16 & 16 & 21,879 \\
\bottomrule\end{tabular}\endgroup\end{center}\end{minipage}

Все 32 численных вызова прошли проверку аргументов и дали правильные итоговые значения с единицами. Среди 11 нечисленных случаев Qwen9 получил 6/11 без tools и 2/11 с tools; Qwen27 --- 9/11 в обеих ветвях. Следовательно, улучшение арифметики не устанавливает надёжность уточнения данных и отказа: у Qwen9 четыре таких случая ухудшились. Ошибок транспорта и усечения по output budget не наблюдалось. Проверке полей и типов JSON соответствуют 24/27 и 26/27 ответов Qwen9, 27/27 в обеих ветвях Qwen27. С дополнительным требованием разрешённого tool ID Qwen9 получает 22/27 и 19/27, Qwen27 --- 27/27 и 27/27; четыре нарушения типов уже неверны по исходному scorer, поэтому отдельная проверка баллы не меняет.


Максимум 512 токенов задан для одного вызова; успешная численная tool-ветвь обычно делает два и получает вдвое больший суммарный бюджет. Обратная связь инструмента меняется вместе с числом вызовов. Указана сумма задержек запросов, а не полная задержка приложения. Это один matched development-опыт на прежнем материале. Пути/размеры весов сверены; полный SHA относится к прежней квитанции и во время inference повторно не рассчитывался.


Сохранены исходные допуски каждого случая: 1e-8 у двенадцати высокоточных численных эталонов, 0,001/0,005/0,01 у четырёх опубликованных округлённых примеров. Отдельная диагностика использует max(исходный допуск, 1e-4) и не заменяет строгие баллы; в исторической серии она не увеличила 1/16 и 6/16. Допуск относится к формуле и инструкции, а не агрономической безопасности. Формульный исполнитель --- технический эталон реализации, не независимое подтверждение правильности полевого решения.


\section*{6. Дополнительные облачные наблюдения}

Mercury2.5 instant/512 на MERA DEV60 дал 31/59 строгих совпадений, включая четыре неуспешные или усечённые попытки. Replay эскалации от Qwen9 с 37/59 не улучшил EM: лучший ненулевой сохранённый порог дал 35/59. Доли обращений симулированы; живые задержка каскада и экономия денег здесь не измерены. Mercury --- облачная reference-служба, её лидерство среди облачных моделей не установлено.


OpenRouter HOME DEV9 содержит по девять завершённых ответов hosted Qwen3.8-27B, Ling3.0-Flash-Sante и Dots3-Note-Preview: 5/9, 4/9 и 4/9 по HOME scorer. Codex CLI на MERA DEV60 дал 36/59 при medium reasoning и 46 нарушениях представления JSON. Claude Code в прежних auth-probes не дал benchmark-ответов. Это сравнения продуктов/стендов; подписка имеет стоимость, а latency включает startup. Hosted имя не означает одинаковые локальные веса.


В историческом HOME150 tools+cloud получил 139/150, локальный каскад --- 133/150 при 20 вместо 100 cloud calls и не выполнил свой критерий качества. В отдельном Jev DEV правила дали 31/36, Jev --- 30/36; все внешние запросы выросли с 23 до 56. Эти серии объясняют важность учёта инструментов и запросов, но не доказывают пользу облачной маршрутизации в основной MERA-работе. Полное покрытие моделей приведено в приложении репозитория, необязательные checkpoints не задерживают статью.


\section*{7. Ограничения и выводы}

DEV повторно использован при настройке и выборе; TELEMETRY40 и сконструированный Study C не являются независимым внешним TEST. Одна машина, разные квантования и поколения, transport/offload, отключённые рассуждения и короткий бюджет ограничивают обобщение. Сохранены нарушения формата, RAM-censored попытки и смены профиля. Строгий scorer воспроизводим локально, production scorer MERA остаётся неизвестным. Экспертное происхождение набора не гарантирует безошибочность ключей и безопасность рекомендаций.


Qwen9 занимает меньше измеренных ресурсов и быстрее отвечает в рассмотренном профиле; Qwen27 имеет более высокое наблюдаемое качество. Это основание для выбора по явно заданному практическому бюджету, без универсального победителя и доказанной эквивалентности. Инструментальная ветвь должна проверяться по аргументам, исполнению и итоговому ответу. Живое исполнение улучшило численные ответы обеих моделей, сохранив ошибки выбора/уточнения у Qwen9 и дополнительный вызов для финализации. Подготовка RU/EN может продолжаться без ответа MERA; отправка требует согласия авторов, метаданных и решения о правах выбранной площадки. Позднейшие открытые модели можно исследовать при том же бюджете, без обещаний будущих названий, дат и качества.


Протоколы, версии scorer, хэши весов/промптов и записи попыток сохранены. Ограниченные материалы MERA не включаются в редакторский пакет. Обучаемый router не тренировался; новых платных API-вызовов, TEST, отправок и внешних писем не было. Codex существенно участвовал в проверке источников, программировании, исполнении, анализе и подготовке текста. Проверка AI не названа человеческой. Авторы --- Антоний Сарычев и Денис Шестернин; окончательные вклады и согласие требуют их подтверждения. Преподаватель и AI автоматически в авторы не включены. Apache-2.0 относится к software; лицензия рукописи и права данных решаются отдельно.


\begin{thebibliography}{99}\fontsize{12}{14}\selectfont
\bibitem{ref1}Chen, L.; Zaharia, M.; Zou, J. \emph{FrugalGPT: How to Use Large Language Models While Reducing Cost and Improving Performance}. arXiv preprint, 2023. \href{https://arxiv.org/abs/2305.05176}{Primary source}.
\bibitem{ref2}Aggarwal, P.; Madaan, A.; Anand, A.; et al. \emph{AutoMix: Automatically Mixing Language Models}. arXiv preprint, 2023. \href{https://arxiv.org/abs/2310.12963}{Primary source}.
\bibitem{ref3}Ding, D.; Mallick, A.; Wang, C.; et al. \emph{Hybrid LLM: Cost-Efficient and Quality-Aware Query Routing}. arXiv preprint, 2024. \href{https://arxiv.org/abs/2404.14618}{Primary source}.
\bibitem{ref4}Ong, I.; Almahairi, A.; Wu, V.; et al. \emph{RouteLLM: Learning to Route LLMs from Preference Data}. ICLR 2025, 2025. \href{https://proceedings.iclr.cc/paper_files/paper/2025/hash/5503a7c69d48a2f86fc00b3dc09de686-Abstract-Conference.html}{Primary source}.
\bibitem{ref5}Dekoninck, J.; Baader, M.; Vechev, M. \emph{A Unified Approach to Routing and Cascading for LLMs}. ICML 2025, PMLR 267:12987-13010, 2025. \href{https://proceedings.mlr.press/v267/dekoninck25a.html}{Primary source}.
\bibitem{ref6}Hu, Q. J.; Bieker, J.; Li, X.; et al. \emph{RouterBench: A Benchmark for Multi-LLM Routing System}. arXiv preprint, 2024. \href{https://arxiv.org/abs/2403.12031}{Primary source}.
\bibitem{ref7}Kossen, J.; Han, J.; Razzak, M.; et al. \emph{Semantic Entropy Probes: Robust and Cheap Hallucination Detection in LLMs}. arXiv preprint, 2024. \href{https://arxiv.org/abs/2406.15927}{Primary source}.
\bibitem{ref8}Wen, B.; Yao, J.; Feng, S.; et al. \emph{Know Your Limits: A Survey of Abstention in Large Language Models}. arXiv preprint, 2024. \href{https://arxiv.org/abs/2407.18418}{Primary source}.
\bibitem{ref9}Yang, X.; Gao, J.; Xue, W.; et al. \emph{PLLaMa: An Open-source Large Language Model for Plant Science}. arXiv preprint, 2024. \href{https://arxiv.org/abs/2401.01600}{Primary source}.
\bibitem{ref10}Yang, B.; Zhang, Y.; Feng, L.; et al. \emph{AgriGPT: a Large Language Model Ecosystem for Agriculture}. arXiv preprint, 2025. \href{https://arxiv.org/abs/2508.08632}{Primary source}.
\bibitem{ref11}Samuel, D. J.; Skarga-Bandurova, I.; Sikolia, D.; et al. \emph{AgroLLM: Connecting Farmers and Agricultural Practices through Large Language Models for Enhanced Knowledge Transfer and Practical Application}. arXiv preprint, 2025. \href{https://arxiv.org/abs/2503.04788}{Primary source}.
\bibitem{ref12}Yang, B.; Zhang, Y.; Chen, Y.; et al. \emph{AgriAgent: Contract-Driven Planning and Capability-Aware Tool Orchestration in Real-World Agriculture}. arXiv preprint, 2026. \href{https://arxiv.org/abs/2601.08308}{Primary source}.
\bibitem{ref13}Ye, Z.; Wen, Y.; Fan, X.; et al. \emph{AgroTools: A Benchmark for Tool-Augmented Multimodal Agents in Agriculture}. arXiv preprint, 2026. \href{https://arxiv.org/abs/2605.22366}{Primary source}.
\bibitem{ref14}Yao, S.; Shinn, N.; Razavi, P.; et al. \emph{tau-bench: A Benchmark for Tool-Agent-User Interaction in Real-World Domains}. arXiv preprint, 2024. \href{https://arxiv.org/abs/2406.12045}{Primary source}.
\bibitem{ref15}Liu, W.; Huang, X.; Zeng, X.; et al. \emph{ToolACE: Winning the Points of LLM Function Calling}. ICLR 2025, 2025. \href{https://proceedings.iclr.cc/paper_files/paper/2025/hash/663865ea167425c6c562cb0b6bcf76c7-Abstract-Conference.html}{Primary source}.
\bibitem{ref16}Ma, S.; Wang, H.; Ma, L.; et al. \emph{The Era of 1-bit LLMs: All Large Language Models are in 1.58 Bits}. arXiv preprint, 2024. \href{https://arxiv.org/abs/2402.17764}{Primary source}.
\bibitem{ref17}Gemma Team; et al. \emph{Gemma 4 Technical Report}. arXiv preprint, 2026. \href{https://arxiv.org/abs/2607.02770}{Primary source}.
\bibitem{ref18}Kpodo, J.; Kordjamshidi, P.; Nejadhashemi, A. P. \emph{AgXQA: A benchmark for advanced Agricultural Extension question answering}. Computers and Electronics in Agriculture 225:109349, 2024. \href{https://doi.org/10.1016/j.compag.2024.109349}{Primary source}.
\bibitem{ref19}Allen, R. G.; Pereira, L. S.; Raes, D.; Smith, M. \emph{Crop evapotranspiration - Guidelines for computing crop water requirements}. FAO Irrigation and Drainage Paper 56, 1998. \href{https://www.fao.org/4/x0490e/x0490e00.htm}{Primary source}.
\end{thebibliography}
\end{document}
